\documentclass{article}

\usepackage[russian]{babel}
\usepackage{fontspec}
\usepackage{geometry}
\usepackage{booktabs}
\usepackage{longtable}
\usepackage{array}
\usepackage{multirow}
\usepackage{enumitem}
\usepackage{titlesec}
\usepackage{xcolor}
\usepackage{siunitx}

\usepackage{pdfpages}

\usepackage{tablefootnote}

% Дополнительные пакеты для данного документа
\usepackage{caption}
\usepackage{subcaption}
\usepackage{listings}
% \usepackage{hyperref}

\geometry{a4paper, margin=25mm}
\setmainfont{Liberation Serif}
\setsansfont{Liberation Sans}

\definecolor{adblue}{RGB}{0,84,166}

\titleformat{\section}{\large\bfseries\color{adblue}}{}{0em}{}[\titlerule]
\titleformat{\subsection}{\normalsize\bfseries}{}{0em}{}
\titleformat{\subsubsection}{\normalsize\itshape}{}{0em}{}

\sisetup{output-decimal-marker={,}}

\usepackage{tikz}
\usetikzlibrary{positioning, calc, backgrounds, math, 3d, perspective, arrows.meta, graphs, graphdrawing, fit, shapes.geometric, trees}
    

\usepackage{pgfplots} % для построения графиков
\pgfplotsset{compat=1.18} % версия для совместимости
\usepackage{tkz-euclide}

% --- Подключение пакета алгоритмов (без флага algo2e, чтобы вернуть \begin{algorithm}) ---
\usepackage[ruled,vlined]{algorithm2e}
%\usepackage{caption}
\newcommand{\Gray}[1]{\textcolor{gray}{#1}} 
% Настройка подписей
\DeclareCaptionLabelFormat{gost}{#1 #2}
\captionsetup[table]{labelformat=gost,labelsep=endash,justification=justified,singlelinecheck=false,font=normal}
\captionsetup[figure]{labelformat=gost,labelsep=endash,justification=centering,font=normal}

% Локализация algorithm2e на русский (ОБЯЗАТЕЛЬНО после загрузки пакета)
\SetAlgorithmName{Алгоритм}{}{}
\SetAlgoCaptionSeparator{---}
\SetKwInput{Input}{Вход}
\SetKwInput{Output}{Выход}

\usepackage{url}
% Говорим пакету использовать моноширинный шрифт (как у texttt)
\Urlmuskip=0mu plus 1mu % гибкие пробелы для формул, если нужно
\DeclareUrlCommand\code{\urlstyle{tt}} 

% Библиотека схемотехники
\usepackage[siunitx, RPvoltages, american]{circuitikzgit}
\ctikzset{
    amplifiers/fill=cyan,
    sources/fill=green!20,
    diodes/fill=white,
    resistors/fill=violet,
    grounds/scale=1.2,
}

\usepackage{iftex}

\ifXeTeX % Если используется TeTeX или TeLaTeX
  \typeout{Режим XeTeX}
  \usepackage{xltxtra}
  \usepackage{fontspec}
  \usepackage{polyglossia}
  \setmainlanguage{russian}
  \setotherlanguage{english}
  \setkeys{russian}{babelshorthands=true} % По идее включает <<, >>, -- ---

  \setmainfont{Times New Roman}[Ligatures=TeX]
  \setromanfont{Times New Roman}[Ligatures=TeX]
  \setsansfont{Arial}[Ligatures=TeX]
  \setmonofont{Courier New}[Ligatures=TeX] 

  \newfontfamily{\cyrillicfont}{Times New Roman}
  \newfontfamily{\cyrillicfontrm}{Times New Roman}
  \newfontfamily{\cyrillicfonttt}{Courier New}
  \newfontfamily{\cyrillicfontsf}{Arial}
\fi

\ifLuaTeX % Если используется LuaLatex
  \typeout{Режим LuaTeX}
%  \usepackage{luatextra}
  \usepackage{fontspec}
  \defaultfontfeatures{Ligatures={TeX}}
  \usepackage{polyglossia}
  \setmainlanguage{russian}
  \setotherlanguage{english}
  \setmainfont{Times New Roman}[Script=Cyrillic, Ligatures=TeX]
  \setromanfont{Times New Roman}[Script=Cyrillic, Ligatures=TeX]
  \setsansfont{Arial}[Script=Cyrillic, Ligatures=TeX]
  \setmonofont{Courier New}[Ligatures=TeX]
  \newfontfamily{\cyrillicfont}{Times New Roman}[Script=Cyrillic, Ligatures=TeX]
  \newfontfamily{\cyrillicfontrm}{Times New Roman}[Script=Cyrillic, Ligatures=TeX]
  \newfontfamily{\cyrillicfonttt}{Courier New}
  \newfontfamily{\cyrillicfontsf}{Arial}[Script=Cyrillic, Ligatures=TeX]

  \usegdlibrary{trees}
  \usegdlibrary{force}
\fi

\ifPDFTeX % Если используется 8битный PDFTex
  \typeout{8-битный режим pdfTex}
  % For pdfTeX we must use old
  % 8-bit font technologies
  \usepackage[T2A]{fontenc}
  \usepackage[english,russian]{babel}
  \usepackage[utf8]{inputenc}
  \usepackage[english, russian]{babel}

  %\usepackage{newtxmath}
  \usepackage{amsmath}
  \usepackage{amsthm}
  \usepackage{amssymb}
  \usepackage{amsfonts}
  \usepackage{mathtext}
\fi

\usepackage{amsthm}  % Возможности AMSMath
\usepackage{amsmath}
\usepackage{amssymb}

\usepackage{esvect} % Для красивых стрелок
\usepackage{bm} % Для жирных символов

%Hyphenation rules
\usepackage{hyphenat}
\hyphenation{ма-те-ма-ти-ка вос-ста-нав-ли-вать ото-бра-жа-ют-ся сге-не-ри-ро-ван-ном сим-во-лы}
% \usepackage[english, russian]{babel}

\usepackage{graphicx}
\graphicspath{ {./images/} }

\usepackage{empheq}

\usepackage{float}

% \usepackage{comment}

\usepackage[hidelinks,
    unicode=true,          % поддерживает Unicode в закладках
    psdextra=true,      % автоматически конвертирует простую математику в текст
    % focusdefault=false
]{hyperref}


% Окружение "Теорема"
\newtheorem{theorem}{Теорема}[section]
\newtheorem{corollary}{Corollary}[theorem]
\newtheorem{lemma}[theorem]{Lemma}

% Окружение "Определение"
\theoremstyle{definition} % Прямой шрифт, нумерация
\newtheorem{definition}{Определение}[section]
\newtheorem{example}{Пример}[section]

% Окружение "Замечание"
\theoremstyle{remark} % Прямой шрифт, без нумерации
% Окуржение "Решение"
\newtheorem*{solution}{Решение}
\newtheorem*{remark}{Замечание}

\usepackage{polynom}
\usepackage{tabularx}
\usepackage{cancel}

\addto\captionsrussian{%
  \renewcommand{\figurename}{Рис.}%
  \renewcommand{\tablename}{Табл.}%
}

\renewcommand{\arraystretch}{1.5} % Высота строки таблицы

\renewcommand{\labelitemi}{\(\circ\)} % Бюллет -- пустой кружок в itemize


\title{Описание алгоритма работы программы восстановления изображений}
\author{}
\date{}

\begin{document}
\maketitle

\section*{Введение}
Программа реализует пайплайн генерации синтетического датасета, обучения нейронной сети NAFNet с комбинированной функцией потерь (L1 + Focal Frequency Loss) и инференса. Ниже представлены ключевые алгоритмы в виде псевдокода.

\section{Генерация обучающей пары (LQ, HQ)}
\label{sec:generation}

\begin{algorithm}[H]
\SetAlgoLined
\KwIn{Исходное резкое RGB-изображение $I_{\text{orig}}$ (uint8, $H\times W\times 3$), конфигурация $\mathcal{C}$}
\KwOut{Пара: $(\text{LQ}_{\text{packed}}, \text{HQ}_{\text{target}})$, где LQ – упакованный RGGB (uint16, $H_{\text{small}}\times W_{\text{small}}\times 4$), HQ – резкое RGB (uint8, $H\times W\times 3$)}
\BlankLine
\SetKwFunction{downscale}{downscale}
\SetKwFunction{psfKernel}{createPSF}
\SetKwFunction{filter2D}{filter2D}
\SetKwFunction{bayerMask}{applyBayerMask}
\SetKwFunction{addNoise}{addNoise}
\SetKwFunction{pack}{packSubchannels}

$\text{HQ}_{\text{target}} \leftarrow I_{\text{orig}}$ \tcc*[r]{сохраняется без изменений}
\If{$\mathcal{C}.downscale > 1$}{
    $I_{\text{small}} \leftarrow \downscale(I_{\text{orig}}, \mathcal{C}.downscale)$ \tcc*[r]{билинейное уменьшение}
}
\Else{
    $I_{\text{small}} \leftarrow I_{\text{orig}}$
}
$I_{\text{small}} \leftarrow I_{\text{small}} / 255.0$ \tcc*[r]{приведение к $[0,1]$, float32}

$\text{kernel} \leftarrow \text{psfKernel}(\sigma=\mathcal{C}.\text{psf\_sigma})$
$I_{\text{blur}} \leftarrow \text{zeros\_similar}(I_{\text{small}})$
\For{$c \in \{R,G,B\}$}{
    $I_{\text{blur}}[:,:,c] \leftarrow \text{filter2D}(I_{\text{small}}[:,:,c], \text{kernel})$ \;
}

$B_{\text{float}} \leftarrow \bayerMask(I_{\text{blur}}, \text{pattern=RGGB})$ \tcc*[r]{одноканальный float [0,1]}

\If{$\mathcal{C}.add\_noise$}{
    $B_{\text{float}} \leftarrow \addNoise(B_{\text{float}}, \mathcal{C}.snr\_db, \mathcal{C}.correlated, \text{kernel})$
}

$B_{16} \leftarrow \text{clip}(B_{\text{float}} \times 65535, 0, 65535).\text{astype(uint16)}$
$\text{LQ}_{\text{packed}} \leftarrow \pack(B_{16})$ \tcc*[r]{$(H_{\text{small}}/2, W_{\text{small}}/2, 4)$}

\KwRet{$\text{LQ}_{\text{packed}}$, $\text{HQ}_{\text{target}}$}
\caption{Генерация синтетической пары (LQ, HQ)}
\end{algorithm}


\section{Один шаг обучения модели}

\begin{algorithm}[H]
\SetAlgoLined
\KwIn{Батч LQ-тензоров $\mathbf{X}$ размерности $(B,4,h,w)$, батч HQ-тензоров $\mathbf{Y}$ размерности $(B,3,2h,2w)$, модель $M$ с параметрами $\theta$, оптимизатор $\mathcal{O}$, функция потерь $\mathcal{L}_{\text{comb}}$, максимальная норма градиентов $g_{\max}$}
\KwOut{Обновлённые параметры $\theta$, значение потерь $L$, метрики (PSNR, L1, FFL)}
\BlankLine
$\mathbf{X}, \mathbf{Y} \leftarrow \text{toDevice}(\mathbf{X}, \mathbf{Y})$
$\mathcal{O}.\text{zero\_grad}()$
$\mathbf{P} \leftarrow M(\mathbf{X})$ \tcc*[r]{прямой проход, $\mathbf{P}$ размерности $(B,3,2h,2w)$}
\eIf{$\mathcal{L}_{\text{comb}}$ — комбинированная потеря}{
    $L_{\text{L1}} \leftarrow \frac{1}{N}\sum |\mathbf{P} - \mathbf{Y}|$
    $L_{\text{FFL}} \leftarrow \mathcal{L}_{\text{FFL}}(\mathbf{P}, \mathbf{Y})$ \tcc*[r]{вычисление через 2D БПФ амплитуд}
    $L \leftarrow w_{\text{L1}}\cdot L_{\text{L1}} + w_{\text{FFL}}\cdot L_{\text{FFL}}$
    $(L_{\text{L1\_val}}, L_{\text{FFL\_val}}) \leftarrow (L_{\text{L1}}, L_{\text{FFL}})$
}{
    $L \leftarrow \frac{1}{N}\sum |\mathbf{P} - \mathbf{Y}|$
    $(L_{\text{L1\_val}}, L_{\text{FFL\_val}}) \leftarrow (L, 0)$
}
$L.\text{backward}()$
$\text{gradNorm} \leftarrow \text{clip\_grad\_norm}(\theta, g_{\max})$ \tcc*[r]{клиппинг нормы градиентов}
$\mathcal{O}.\text{step}()$
Вычислить PSNR: $\text{PSNR} = 20\log_{10}\left(\frac{1}{\sqrt{\text{MSE}(\mathbf{P}, \mathbf{Y})}}\right)$
\KwRet{$L$, $\text{PSNR}$, $L_{\text{L1\_val}}$, $L_{\text{FFL\_val}}$}
\caption{Один итерационный шаг обучения}
\end{algorithm}


\section{Функция Focal Frequency Loss (FFL)}

\begin{algorithm}[H]
\SetAlgoLined
\KwIn{Prediction $\mathbf{P} \in \mathbb{R}^{B\times C\times H\times W}$, target $\mathbf{Y} \in \mathbb{R}^{B\times C\times H\times W}$, parameters $\alpha \ge 0$, $w_{\text{loss}}$}
\KwOut{FFL value}
\KwOut{Значение FFL}
\BlankLine
$\mathcal{F}_P \leftarrow \text{fft2}(\mathbf{P})$ \tcc*[r]{комплексное 2D БПФ}
$\mathcal{F}_Y \leftarrow \text{fft2}(\mathbf{Y})$
$\mathcal{F}_P \leftarrow \text{fftshift}(\mathcal{F}_P)$
$\mathcal{F}_Y \leftarrow \text{fftshift}(\mathcal{F}_Y)$
$A_P \leftarrow |\mathcal{F}_P|$ \tcc*[r]{амплитудный спектр}
$A_Y \leftarrow |\mathcal{F}_Y|$
$D \leftarrow (A_P - A_Y)^2$
$m \leftarrow \max(D) + \epsilon$
$W \leftarrow (D / m)^{\alpha}$ \tcc*[r]{фокальный вес}
$\text{loss} \leftarrow \text{mean}(W \cdot D) \times w_{\text{loss}}$
\KwRet{loss}
\caption{Focal Frequency Loss}
\end{algorithm}

\section{Псевдокод загрузки датасета с аугментациями}

\begin{algorithm}[H]
\SetAlgoLined
\KwIn{Путь к папкам LQ и HQ, размеры $h_{\text{lq}}, w_{\text{lq}}$, флаги $\text{useFlip},\text{useRot}$}
\KwOut{Тензоры $(\mathbf{X}, \mathbf{Y})$}
\BlankLine
Загрузить LQ: $\text{lq} \leftarrow \text{imread}(path_{\text{lq}}) / 65535.0$ \tcc*[r]{$(H_l,W_l,4)$}
Загрузить HQ: $\text{hq} \leftarrow \text{imread}(path_{\text{hq}}) / 255.0$ \tcc*[r]{$(H_h,W_h,3)$}
\If{$(H_h,W_h) \neq (2H_l,2W_l)$}{
    $\text{hq} \leftarrow \text{resize}(\text{hq}, (2H_l,2W_l))$ \tcc*[r]{защита от несоответствия}
}
Выбрать случайные координаты $t_l \in [0, H_l-h_{\text{lq}}]$, $l_l \in [0, W_l-w_{\text{lq}}]$
$\text{lq\_crop} \leftarrow \text{lq}[t_l:t_l+h_{\text{lq}}, l_l:l_l+w_{\text{lq}}, :]$
$\text{hq\_crop} \leftarrow \text{hq}[2t_l:2t_l+2h_{\text{lq}}, 2l_l:2l_l+2w_{\text{lq}}, :]$
$\mathbf{X} \leftarrow \text{tensor}(\text{lq\_crop.transpose(2,0,1)})$, $\mathbf{Y} \leftarrow \text{tensor}(\text{hq\_crop.transpose(2,0,1)})$

\uIf{$\text{useFlip}$}{
    \If{random() $>0.5$}{$\mathbf{X},\mathbf{Y} \leftarrow \text{hflip}(\mathbf{X}), \text{hflip}(\mathbf{Y})$}
    \If{random() $>0.5$}{$\mathbf{X},\mathbf{Y} \leftarrow \text{vflip}(\mathbf{X}), \text{vflip}(\mathbf{Y})$}
}
\If{$\text{useRot}$}{
    $k \leftarrow \text{randint}(0,3)$
    \If{$k>0$}{$\mathbf{X},\mathbf{Y} \leftarrow \text{rotate}(\mathbf{X}, k\cdot 90^\circ), \text{rotate}(\mathbf{Y}, k\cdot 90^\circ)$}
}
\KwRet{$(\mathbf{X}, \mathbf{Y})$}
\caption{Получение элемента датасета с синхронными аугментациями}
\end{algorithm}

\end{document}
